\section{Introduction}
\label{sec:intro}

A security operator watches a live camera feed and asks: \emph{``Did anyone enter the building in the last five minutes?''}
Someone wearing smart glasses asks: \emph{``What did I put in my bag?''}
Both questions target a video stream that is still running. Standard video question answering (VideoQA) cannot help here---it assumes the full video is available before inference starts.

Recent models like Video-ChatGPT~\cite{maaz2023videochatgpt}, VideoLLaVA~\cite{lin2023videollava}, and LLaVA-Next-Video~\cite{zhang2024llavanextvideo} share one assumption: the video is finite and fully available at query time.
Frame sampling, global temporal pooling, and bidirectional attention all depend on that assumption.
A live stream has no end. Feed one into these models and three problems surface:

\begin{enumerate}[leftmargin=*,itemsep=2pt]
    \item \textbf{Unbounded memory.} Encoding every frame from a continuous stream means memory and compute grow linearly without bound.
    \item \textbf{Temporal ambiguity.} \emph{``What is happening?''} targets the present; \emph{``Was there a dog earlier?''} targets the past. Without explicit temporal reasoning, the model mixes these contexts.
    \item \textbf{Latency.} Re-processing the full stream for each query adds delay that precludes real-time interaction.
\end{enumerate}

\textbf{\ours} solves all three. It is a streaming vision-language system for interactive Q\&A on live video, built from three components:

\begin{enumerate}[leftmargin=*,itemsep=2pt]
    \item \textbf{Semantic Keyframe Memory (SKM).}
    A fixed bank of $N$ entries stores the most informative frames, scored by visual novelty and temporal coverage.

    \item \textbf{Temporal Query Router (TQR).}
    Classifies each question into one of three scopes---\emph{instantaneous}, \emph{recent}, or \emph{historical}---and passes only the matching memory slice forward.

    \item \textbf{Stream-Fused Generator (SFG).}
    BLIP~\cite{li2022blip} runs captioning and VQA on filtered keyframes; Flan-T5~\cite{chung2022flanT5} fuses those observations into a single response (254\,ms average on A100 GPU).
\end{enumerate}

We introduce \emph{LiveQA-Bench}, spanning 55 diverse video streams with explicit scope labels---the first VideoQA benchmark that enforces causal access with temporal scope annotations.
Per-scope results expose clear patterns: historical questions score highest (77.1\%) because the full SKM provides rich context; instant (46.2\%) depends on single-frame perception; recent (43.2\%) requires localizing context within a 30\,s window.
Ablations validate every component: removing TQR drops accuracy 3.9 points (instant collapses 15.3 points), replacing SKM with FIFO costs 1.2 points, and compact memory ($N\!=\!16$) achieves the best accuracy (57.2\%).
Published results on OVO-Bench~\cite{li2025ovobench} confirm a large gap between offline and streaming models; our architecture targets this gap with importance-scored memory and temporal routing.

To summarize our contributions:
\begin{enumerate}[leftmargin=*,itemsep=1pt]
  \item We formalize the \emph{streaming VideoQA} task with a causal access constraint and temporal scope taxonomy, and contribute LiveQA-Bench with per-scope annotations.
  \item We propose \ours, combining importance-scored memory, temporal routing, and a two-stage perception-synthesis pipeline, all using frozen pre-trained checkpoints.
  \item Per-scope error analysis identifies specific failure modes---scope misclassification, sparse keyframes, generic generation---and points to concrete improvements.
\end{enumerate}

A live web demo lets users ask questions about their own webcam feed in real time.
